% Dataset overview table. Every cell and caption claim verified against the
% codebase: input dims / label spaces from results.json (num_features,
% num_classes), fractions and splits from the loaders and configs, one-hot /
% MinMax fit scope from src/data/*.py, subsample implementation (seeded
% uniform random, NOT stratified) from the three subsampling loaders, and the
% empty-Probe reason from src/data/cseids2018.py (no port-scan attacks exist
% in the dataset). No em dashes. Requires: \usepackage{booktabs}
\begin{table*}[!ht]
\centering
\renewcommand{\arraystretch}{1.1}
\setlength{\tabcolsep}{4pt}
\caption{Overview of the general-purpose and IoT/IIoT/IoMT network intrusion
detection datasets used in this study. \textbf{Input dim.} is the model-input
width after dataset-specific cleaning and preprocessing, including categorical
one-hot encoding, selected log1p transformations, and MinMax scaling to
$[0,1]$; targets, identifiers, and split metadata are excluded. Numeric
scaling is fit using the training partition only and is then applied unchanged
to the validation and test partitions. To maintain a fixed input schema, the
one-hot category vocabulary is constructed from the union of the train,
validation, and test partitions. This transductive schema-definition step uses
held-out category identities, but not labels, category frequencies, or numeric
summary statistics, and is applied identically to all models. \textbf{Native labels} describes the label granularity in the
flow records consumed by each loader; capture days, capture scenarios, and
broader attack families are not counted as classes. \textbf{Label space} is
the configured classifier output space. For datasets with provider-defined
train/test partitions, the provider test set is retained and 10\% of the
available training partition is held out for validation. For datasets without
provider-defined partitions, a class-stratified train/test split is created
first and 10\% of the resulting training portion is then held out for
validation; the 70/30 and 80/20 splits therefore correspond to effective
63/7/30 and 72/8/20 train/validation/test partitions. \textbf{Frac.} is the
seeded uniform-random fraction of flows retained before partitioning, except
where otherwise stated. The SNN students and neural
baselines are trained using batch size 128, Adam, 10 epochs, and five seeds
(42--46); the GBDT teacher uses a separate gradient-boosting configuration.
$^{a}$CTU-13 comprises 13 capture scenarios containing Botnet, Normal, and
Background traffic; Background flows are excluded per the canonical protocol,
yielding a binary Normal-versus-Botnet task.
$^{b}$ISCXIDS2012 is a seven-day capture whose processed flow records use
binary Normal/Attack labels.
$^{c}$KDDCup99 and NSL-KDD retain fine-grained attack names in the source
files; their loaders apply the standard \{Normal, DoS, Probe, R2L, U2R\}
grouping before model construction, including the mapping of attack types
that occur only in the test partitions.
$\ddagger$For CIC-IDS2017 and CSE-CIC-IDS2018 only, native fine-grained
labels are mapped to the KDD five-class taxonomy \{Normal, DoS, Probe, R2L,
U2R\} to align these experiments with the KDDCup99 and NSL-KDD tasks; the
complete mapping and post-mapping class counts are reported in the appendix.
$*$CSE-CIC-IDS2018 contains no port-scan attack scenario. Under the mapping
used in this study, port-scanning is the source of the Probe class; Probe is
therefore empty by construction, and evaluation covers the four populated
classes of the configured five-class label space.
$\dagger$For CICIoMT2024, a 25\% seeded uniform-random sample is drawn
separately from the provider training and test partitions, after which 10\%
of the sampled training data is held out for validation.}
\label{tab:datasets_overview}
{\footnotesize
\begin{tabular}{l c c c l c l}
\toprule
\textbf{Dataset} & \textbf{Input} & \textbf{Native} & \textbf{Label} &
\textbf{Partition protocol} & \textbf{Frac.} & \textbf{Source} \\
& \textbf{dim.} & \textbf{labels} & \textbf{space} & & & \\
\midrule
\multicolumn{7}{l}{\textit{General-purpose network IDS datasets}} \\
\midrule

KDDCup99
& 119 & Fine-grained$^{c}$ & 5
& Official 10\% train / corrected test; 90/10 val
& 1.0
& ACM SIGKDD, 1999 \cite{kddcup1999} \\

NSL-KDD
& 122 & Fine-grained$^{c}$ & 5
& Official KDDTrain+ / KDDTest+; 90/10 val
& 1.0
& Tavallaee et al., \textit{CISDA} 2009 \cite{tavallaee2009detailed} \\

CTU-13
& 37 & 3$^{a}$ & 2
& Stratified 70/30; 90/10 val
& 1.0
& Garc\'ia et al., \textit{Comput.\,\&\,Secur.} 2014 \cite{garcia2014empirical} \\

ISCXIDS2012
& 175 & 2$^{b}$ & 2
& Stratified 70/30; 90/10 val
& 1.0
& Shiravi et al., \textit{Comput.\,\&\,Secur.} 2012 \cite{shiravi2012toward} \\

CIC-IDS2017$^{\ddagger}$
& 78 & 15 & 5
& Stratified 70/30; 90/10 val
& 1.0
& Sharafaldin et al., \textit{ICISSP} 2018 \cite{sharafaldin2018toward} \\

CSE-CIC-IDS2018$^{\ddagger *}$
& 78 & 15 & 5   % 82 before the leakage repair dropped Flow ID, Src IP, Src Port, Dst IP
& Stratified 70/30; 90/10 val
& 0.5
& CSE-CIC, 2018 \cite{karatas2020increasing} \\

UNSW-NB15
& 196 & 10 & 10
& Official 175{,}341 / 82{,}332; 90/10 val
& 1.0
& Moustafa and Slay, \textit{MilCIS} 2015 \cite{moustafa2015unsw} \\

\midrule
\multicolumn{7}{l}{\textit{IoT, IIoT, and IoMT IDS datasets}} \\
\midrule

CICIoT2023
& 46 & 34 & 8
& Stratified 80/20; 90/10 val
& 0.05
& Neto et al., \textit{Sensors} 2023 \cite{neto2023ciciot2023} \\

CICIoMT2024$^{\dagger}$
& 45 & 19 & 6
& Provider train/test partition; 90/10 val
& 0.25
& Dadkhah et al., \textit{Internet of Things} 2024 \cite{dadkhah2024ciciomt2024} \\

Edge-IIoTset
& 64 & 15 & 14  % 95 / 15 before the repair (zero-spelling canonicalisation, two fields dropped, MITM removed)
& Stratified 80/20; 90/10 val
& 1.0
& Ferrag et al., \textit{IEEE Access} 2022 \cite{ferrag2022edge} \\

TON-IoT (Network)
& 115 & 10 & 10
& Stratified 80/20; 90/10 val
& 1.0
& Moustafa, \textit{Sustain.\ Cities Soc.} 2021 \cite{moustafa2021new} \\

\bottomrule
\end{tabular}}
\end{table*}
